% Source: https://artofproblemsolving.com/community/c3h2267005_probability_question
% Stable identity: aops-topic-2267005-q01
Sam does not know the answers to $5$ questions of a multiple-choice test; each has $5$ choices. For each of these questions he can either
\begin{itemize}
\item guess at random, and then he is right with probability $\dfrac{1}{5}$; or
\item reason it out, and then he is right exactly when his reasoning is correct, which happens with probability $p$ (where $0\le p\le1$).
\end{itemize}
The results on different questions are independent. Sam wants to maximise the probability of getting all $5$ questions right.
\begin{enumerate}
\item[(a)] If he reasons on $k$ of the questions and guesses on the other $5-k$, what is the probability that he gets all $5$ right?
\item[(b)] For which values of $p$ should he reason on all five questions, and for which should he guess on all five? Can a combination of the two methods ever give a strictly better chance than both pure strategies?
\end{enumerate}
